% ============================================================
% preamble.tex —— 导言区（Lead 独占维护；译者只读，不得修改）
% 编译：xelatex ×2，必须先把 TMP/TEMP 指到 work/tmp
% ============================================================
\usepackage[a4paper,top=2.6cm,bottom=2.6cm,left=2.6cm,right=2.6cm]{geometry}
\usepackage{amsmath,amssymb,amsthm,mathtools}
\usepackage{graphicx}
\usepackage{xcolor}
\usepackage{tikz}
\usepackage{enumitem}
\usepackage{array,booktabs,longtable,multirow,tabularx}
\usepackage{fancyhdr}
\usepackage[colorlinks=true,linkcolor=blue!55!black,citecolor=blue!55!black,urlcolor=blue!55!black]{hyperref}

% ---------- 字体（fontset=none + 显式文件；不要用 fontconfig） ----------
\setCJKmainfont[Path=C:/Windows/Fonts/,BoldFont=simhei.ttf,ItalicFont=simkai.ttf]{simsun.ttc}
\setCJKsansfont[Path=C:/Windows/Fonts/]{simhei.ttf}
\setCJKmonofont[Path=C:/Windows/Fonts/]{simfang.ttf}

% ---------- 排版细节 ----------
\linespread{1.15}
\setlength{\emergencystretch}{3em}
\setcounter{tocdepth}{2}
\setcounter{secnumdepth}{3}
\renewcommand{\proofname}{证明}
\newenvironment{proofof}[1]{\begin{proof}[#1]}{\end{proof}}

% 书签里禁止数学（防 hyperref 警告）：不把数学写进 \chapter/\section 标题即可

% ---------- 定理类环境（编号一律手工给定：\begin{theorem}[1.1]） ----------
\theoremstyle{plain}
\newtheorem*{theorem}{定理}
\newtheorem*{proposition}{命题}
\newtheorem*{lemma}{引理}
\newtheorem*{corollary}{推论}
\newtheorem*{claim}{断言}
\newtheorem*{fact}{事实}
\newtheorem*{conjecture}{猜想}
\newtheorem*{hypothesis}{假设}
\theoremstyle{definition}
\newtheorem*{definition}{定义}
\newtheorem*{example}{例}
\newtheorem*{exercise}{习题}
\newtheorem*{construction}{构造}
\newtheorem*{setup}{设定}
\newtheorem*{algorithm}{算法}
\newtheorem*{convention}{约定}
\newtheorem*{notation}{记号}
\newtheorem*{question}{问题}
\newtheorem*{problem}{问题}
\theoremstyle{remark}
\newtheorem*{remark}{注}
\newtheorem*{warning}{警告}

% ---------- 数学宏（全局统一；译者不得自造 \newcommand） ----------
\newcommand{\R}{\mathbb{R}}
\newcommand{\C}{\mathbb{C}}
\newcommand{\Z}{\mathbb{Z}}
\newcommand{\Q}{\mathbb{Q}}
\newcommand{\N}{\mathbb{N}}
\newcommand{\PP}{\mathbb{P}}
\newcommand{\HH}{\mathbb{H}}
\newcommand{\DD}{\mathbb{D}}
\newcommand{\OO}{\mathcal{O}}
\newcommand{\g}{\mathfrak{g}}
\newcommand{\h}{\mathfrak{h}}
\newcommand{\n}{\mathfrak{n}}
\renewcommand{\sl}{\mathfrak{sl}}
\newcommand{\gl}{\mathfrak{gl}}
\newcommand{\so}{\mathfrak{so}}
\newcommand{\su}{\mathfrak{su}}
\renewcommand{\sp}{\mathfrak{sp}}
\newcommand{\Tr}{\operatorname{Tr}}
\newcommand{\tr}{\operatorname{tr}}
\newcommand{\rank}{\operatorname{rank}}
\newcommand{\diag}{\operatorname{diag}}
\newcommand{\Spec}{\operatorname{Spec}}
\newcommand{\Hom}{\operatorname{Hom}}
\newcommand{\End}{\operatorname{End}}
\newcommand{\Aut}{\operatorname{Aut}}
\newcommand{\GL}{\operatorname{GL}}
\newcommand{\SL}{\operatorname{SL}}
\newcommand{\PSL}{\operatorname{PSL}}
\newcommand{\PGL}{\operatorname{PGL}}
\newcommand{\SO}{\operatorname{SO}}
\newcommand{\SU}{\operatorname{SU}}
\newcommand{\U}{\operatorname{U}}
\newcommand{\Mat}{\operatorname{Mat}}
\newcommand{\Rep}{\operatorname{Rep}}
\newcommand{\Loc}{\operatorname{Loc}}
\newcommand{\Fib}{\operatorname{Fib}}
\newcommand{\Iso}{\operatorname{Iso}}
\newcommand{\im}{\operatorname{im}}
\newcommand{\coker}{\operatorname{coker}}
\newcommand{\sgn}{\operatorname{sgn}}
\newcommand{\const}{\operatorname{const}}
\newcommand{\de}{\mathrm{d}}          % 微分 d（正体）
\newcommand{\ee}{\mathrm{e}}
\newcommand{\ii}{\mathrm{i}}
\newcommand{\dd}{\mathrm{d}}
\newcommand{\vol}{\operatorname{vol}}
\newcommand{\area}{\operatorname{area}}
\newcommand{\length}{\operatorname{length}}
\newcommand{\op}[1]{\operatorname{#1}}
\newcommand{\Sp}{\operatorname{Sp}}
\newcommand{\Stab}{\operatorname{Stab}}
\newcommand{\Sym}{\operatorname{Sym}}
\newcommand{\Homeo}{\operatorname{Homeo}}
\newcommand{\tq}{\text{；}}           % 数学环境里的中文分号（需自己包 \text{}）

% ---------- 插图（Lead 已裁好图；译者只引用，不得自造） ----------
% 用法：\Fig[0.75\linewidth]{figures/fig_1_3.png}{1.3}{中文图注}
\newcommand{\Fig}[4][0.75\linewidth]{%
  \begin{figure}[htbp]\centering
    \includegraphics[width=#1]{#2}\par\vspace{3pt}
    {\small 图~#3\quad #4}
  \end{figure}}
\newcommand{\FigNoNum}[3][0.75\linewidth]{%
  \begin{figure}[htbp]\centering
    \includegraphics[width=#1]{#2}\par\vspace{3pt}
    {\small #3}
  \end{figure}}

% ---------- 表格：按原文重排（不造新表） ----------
\renewcommand{\arraystretch}{1.15}

% ---------- 章首摘要（Springer 体例每章有 Abstract） ----------
\newcommand{\ChapterAbstract}[1]{\begin{quote}\small\noindent\textbf{摘要}\quad #1\end{quote}}
% ---------- 定理头格式：去掉 amsthm 自动加的 ASCII 括号 ----------
% 目标：\begin{theorem}[6.3] → 「定理 6.3」；\begin{theorem}[7.2（庞加莱单值化）] → 「定理 7.2（庞加莱单值化）」
\makeatletter
\renewcommand{\thmhead}[3]{\thmname{#1}\thmnumber{ #2}\thmnote{ #3}}
\makeatother
% ---------- 追加环境与宏（2026-11-03 第二批） ----------
\newcommand{\Lag}{\operatorname{Lag}}
\newcommand{\Conf}{\operatorname{Conf}}\newcommand{\PSU}{\operatorname{PSU}}
\newcommand{\USp}{\operatorname{USp}}
\newcommand{\PU}{\operatorname{PU}}
\newcommand{\Spin}{\operatorname{Spin}}
\newcommand{\Pin}{\operatorname{Pin}}
